\chapter{Resultados Complementares}\label{ap:resultados-complementares}

Este apêndice reúne as tabelas e figuras que aprofundam as análises do Capítulo~\ref{cap:resultados}, concentrando-se nos desdobramentos que não foram detalhados no corpo principal do texto: cobertura complementar entre linguagens, distribuição de vereditos por modelo, telemetria de tamanho de código e custo, mapeamento nominal das questões com transição de veredito e auditoria dos erros de extração.

\section{Detalhamento Complementar da PP1: Desempenho por Linguagem}\label{ap:pp1-detalhes}

\subsection{Taxa de Acerto (\texorpdfstring{\passat{1}}{Pass@1}) e Cobertura Complementar}\label{sec:res-rq1-pass1}

A Tabela~\ref{tab:res-pp1-pass1} apresenta os números absolutos e percentuais de aprovação (\veredito{AC} de 125) de cada modelo nas cinco linguagens avaliadas, servindo de referência numérica para a Figura~\ref{fig:res-pp1-ac-linguagem} discutida na Seção~\ref{sec:res-rq1}.

\begin{table}[!htbp]
\centering
\caption[Taxa de acerto Pass@1 por modelo e linguagem na PP1]{Quantidade absoluta de aprovações (\veredito{AC} de 125) e taxa $\passat{1}$ (\%) por modelo e linguagem de programação ($T = 0{,}0$).}
\label{tab:res-pp1-pass1}
\small
\resizebox{\textwidth}{!}{%
\begin{tabular}{lccccccc}
\toprule
Modelo & Python & C & C++ & Java & JavaScript & Média ($\passat{1}$) & Total AC (625) \\
\midrule
DeepSeek v4.1 Flash & \tsecond{118} (94,40\%) & \tfirst{118} (94,40\%) & \tfirst{121} (96,80\%) & \tsecond{117} (93,60\%) & \tsecond{115} (92,00\%) & \tsecond{117,8} (94,24\%) & \tsecond{589} \\
Gemini 3.8 Flash    & \tfirst{121} (96,80\%)  & \tsecond{117} (93,60\%) & \tsecond{119} (95,20\%) & 114 (91,20\%) & \tfirst{119} (95,20\%)  & \tfirst{118,0} (94,40\%)  & \tfirst{590}  \\
GPT-6 Luna          & 109 (87,20\%)           & 110 (88,00\%)           & 113 (90,40\%)           & \tfirst{119} (95,20\%)  & \tsecond{115} (92,00\%) & 113,2 (90,56\%)           & 566           \\
\midrule
\rowcolor{linhaDestaque}\textbf{Agregado (de 375)} & 348 (92,80\%) & 345 (92,00\%) & \tfirst{353} (94,13\%) & \tsecond{350} (93,33\%) & 349 (93,07\%) & -- & 1.745 (93,07\%) \\
\bottomrule
\end{tabular}%
}
\par\smallskip{\footnotesize \tfirst{Vermelho} e \tsecond{azul}: melhor e segundo melhor valor por coluna.}
\end{table}

Além das medidas de dispersão apresentadas na Tabela~\ref{tab:res-pp1-estabilidade} do Capítulo~\ref{cap:resultados}, um aspecto relevante é verificar se os modelos erram sempre as mesmas questões ou se as falhas mudam conforme a linguagem escolhida. A Tabela~\ref{tab:ap-pp1-estabilidade} mostra a cobertura complementar de cada LLM sobre o conjunto de 125 problemas. Embora o número de questões resolvidas de forma unânime nas cinco linguagens (5/5) varie entre 92 (no \textbf{GPT-6 Luna}) e 105 (no \textbf{DeepSeek v4.1 Flash} e no \textbf{Gemini 3.8 Flash}), todos os modelos conseguem resolver entre 123 e 124 questões ($98{,}4\%$ a $99{,}2\%$) quando se considera o acerto em ao menos uma linguagem ($\ge 1$). Esse comportamento reforça que a grande maioria dos erros pontuais na PP1 não decorre da falta de capacidade lógica para resolver o problema, mas de tropeços específicos na sintaxe ou nas particularidades de implementação de cada linguagem.

\begin{table}[!htbp]
\centering
\caption[Cobertura complementar entre linguagens na PP1]{Cobertura complementar de resolução de cada modelo ao longo das cinco linguagens sobre as 125 questões da OBI Sênior.}
\label{tab:ap-pp1-estabilidade}
\small
\resizebox{\textwidth}{!}{%
\begin{tabular}{lccc}
\toprule
Modelo & Acertos nas 5 Linguagens (5/5) & Acerto em ao menos 1 ($\ge 1$) & Erro nas 5 Linguagens (0/5) \\
\midrule
DeepSeek v4.1 Flash & \tfirst{105/125} (84,00\%) & 123/125 (98,40\%)          & 2/125 (1,60\%) \\
Gemini 3.8 Flash    & \tfirst{105/125} (84,00\%) & \tfirst{124/125} (99,20\%) & \tfirst{1/125} (0,80\%) \\
GPT-6 Luna          & 92/125 (73,60\%)           & \tfirst{124/125} (99,20\%) & \tfirst{1/125} (0,80\%) \\
\bottomrule
\end{tabular}%
}
\end{table}

\subsection{Distribuição de Vereditos por Modelo e Linguagem}\label{sec:res-rq1-falhas}

Enquanto o Capítulo~\ref{cap:resultados} (Tabela~\ref{tab:res-pp1-erros-lang}) sintetizou os erros por linguagem de programação, a Tabela~\ref{tab:res-pp1-vereditos-llm} e a Figura~\ref{fig:res-pp1-erros-linguagem} mostram como esses vereditos se dividem entre os três LLMs nas 625 submissões de cada modelo.

\begin{table}[!htbp]
\centering
\caption[Distribuição geral dos vereditos por LLM na PP1]{Distribuição geral dos vereditos do \textit{JudgeService} por LLM nas 625 execuções da PP1 ($T = 0{,}0$).}
\label{tab:res-pp1-vereditos-llm}
\small
\resizebox{\textwidth}{!}{%
\begin{tabular}{lccccccc}
\toprule
Modelo & \chip{corAC}\veredito{AC} & \chip{corWA}\veredito{WA} & \chip{corRE}\veredito{RE} & \chip{corTLE}\veredito{TLE} & \chip{corMLE}\veredito{MLE} & \chip{corCE}\veredito{CE} & Total \\
\midrule
DeepSeek v4.1 Flash & 589 (94,2\%) & 23 (3,7\%) & 1 (0,2\%) & 1 (0,2\%) & 0 (0,0\%) & 11 (1,8\%) & 625 \\
Gemini 3.8 Flash    & 590 (94,4\%) & 24 (3,8\%) & 4 (0,6\%) & 0 (0,0\%) & 2 (0,3\%) & 5 (0,8\%)  & 625 \\
GPT-6 Luna          & 566 (90,6\%) & 47 (7,5\%) & 2 (0,3\%) & 3 (0,5\%) & 1 (0,2\%) & 6 (1,0\%)  & 625 \\
\bottomrule
\end{tabular}%
}
\end{table}

\begin{figure}[!htbp]
\centering
% Gráfico de barras empilhadas: distribuição de erros (Não-AC) por linguagem (Cap. 5, Seção 5.2.2)
\begin{tikzpicture}
\begin{axis}[
  ybar stacked,
  width=0.86\textwidth,
  height=6.2cm,
  bar width=26pt,
  ymin=0, ymax=36,
  ylabel={Quantidade de erros (de 375 execuções)},
  xlabel={Linguagem de programação},
  symbolic x coords={Python,C,C++,Java,JavaScript},
  xtick=data,
  enlarge x limits=0.16,
  ymajorgrids=true,
  grid style={black!12},
  axis lines*=left,
  tick label style={font=\footnotesize},
  label style={font=\small},
  legend style={
    at={(0.5,1.03)},
    anchor=south,
    legend columns=5,
    draw=black!30,
    rounded corners=2pt,
    font=\footnotesize,
    /tikz/every even column/.append style={column sep=6pt}
  }
]
\addplot[fill=corWA, draw=corWA!75!black, line width=0.5pt] coordinates {
  (Python,21) (C,17) (C++,19) (Java,17) (JavaScript,20)
};
\addplot[fill=corCE, draw=corCE!75!black, line width=0.5pt] coordinates {
  (Python,2) (C,13) (C++,2) (Java,3) (JavaScript,2)
};
\addplot[fill=corRE, draw=corRE!75!black, line width=0.5pt] coordinates {
  (Python,1) (C,0) (C++,0) (Java,4) (JavaScript,2)
};
\addplot[fill=corTLE, draw=corTLE!75!black, line width=0.5pt] coordinates {
  (Python,3) (C,0) (C++,1) (Java,0) (JavaScript,0)
};
\addplot[fill=corMLE, draw=corMLE!75!black, line width=0.5pt] coordinates {
  (Python,0) (C,0) (C++,0) (Java,1) (JavaScript,2)
};
\legend{\veredito{WA}, \veredito{CE}, \veredito{RE}, \veredito{TLE}, \veredito{MLE}}

% Totais agregados acima de cada barra empilhada
\node[font=\footnotesize\bfseries, anchor=south] at (axis cs:Python,27) {27};
\node[font=\footnotesize\bfseries, anchor=south] at (axis cs:C,30) {30};
\node[font=\footnotesize\bfseries, anchor=south] at (axis cs:C++,22) {22};
\node[font=\footnotesize\bfseries, anchor=south] at (axis cs:Java,25) {25};
\node[font=\footnotesize\bfseries, anchor=south] at (axis cs:JavaScript,26) {26};
\end{axis}
\end{tikzpicture}

\caption[Distribuição agregada de erros por linguagem na PP1]{Distribuição agregada dos erros (não-\veredito{AC}) por linguagem de programação nas 375 execuções de cada linguagem.}
\label{fig:res-pp1-erros-linguagem}
\end{figure}
\FloatBarrier

Cruzando os vereditos por modelo e linguagem, notam-se comportamentos bem característicos. A linguagem \linguagem{C} concentrou 13 dos 22 erros de compilação (\veredito{CE}) de toda a bateria --- incluindo todos os 5 \veredito{CE}s registrados pelo \textbf{Gemini 3.8 Flash} ---, devido à exigência de inclusão explícita de cabeçalhos e manipulação manual de memória. Já no \textbf{DeepSeek v4.1 Flash}, a inspeção dos arquivos de saída revelou que os seus 11 \veredito{CE}s (\linguagem{Python}: 2; \linguagem{C}: 4; \linguagem{C++}: 1; \linguagem{Java}: 2; \linguagem{JavaScript}: 2) ocorreram porque o modelo atingiu o teto de $131.072$ \textit{tokens} ainda durante o raciocínio interno, encerrando a resposta antes de emitir o bloco de código. Por fim, \linguagem{Python} concentrou 3 dos 4 casos de \veredito{TLE} por conta do custo de interpretação em laços extensos, ao passo que \linguagem{Java} reuniu 4 dos 7 erros em tempo de execução (\veredito{RE}), geralmente ligados à leitura de entrada ou estouro de limites de vetores e pilha.

\subsection{Extensão do Código, Consumo de \textit{Tokens}, Custo e Tempo de CPU}\label{sec:res-rq1-eficiencia}

A Tabela~\ref{tab:res-pp1-eficiencia} detalha as médias de linhas de código (LOC), número de caracteres, \textit{tokens} de saída, custo acumulado em dólares e tempo médio de CPU das soluções aprovadas para cada par modelo--linguagem (dados que originaram a Figura~\ref{fig:res-pp1-loc-tempo} do Capítulo~\ref{cap:resultados}).

\begin{table}[!htbp]
\centering
\caption[Métricas de código, tokens, custo e tempo de CPU na PP1]{Extensão média do código (LOC e caracteres), \textit{tokens} médios de saída, custo total (USD) e tempo médio de CPU nas soluções \veredito{AC} por modelo e linguagem.}
\label{tab:res-pp1-eficiencia}
\small
\resizebox{\textwidth}{!}{%
\begin{tabular}{llrrrrr}
\toprule
Modelo & Linguagem & LOC Médio & Caracteres Médio & \textit{Tokens} Saída Médio & Custo Total (USD) & Tempo Médio \veredito{AC} (s) \\
\midrule
DeepSeek v4.1 Flash & Python     & 42,6 & 911,2   & 8.688,1  & \$0,6519 & 0,2776 \\
DeepSeek v4.1 Flash & C          & 58,7 & 1.259,8 & 11.387,3 & \$0,8544 & 0,1569 \\
DeepSeek v4.1 Flash & C++        & 49,9 & 1.073,1 & 7.541,5  & \$0,5659 & 0,1856 \\
DeepSeek v4.1 Flash & Java       & 60,8 & 1.684,8 & 8.939,3  & \$0,6708 & 0,1365 \\
DeepSeek v4.1 Flash & JavaScript & 46,8 & 1.027,4 & 9.198,4  & \$0,6902 & 0,0929 \\
\midrule
Gemini 3.8 Flash    & Python     & 47,6 & 1.151,9 & 435,6    & \$0,2828 & 0,2973 \\
Gemini 3.8 Flash    & C          & 60,9 & 1.334,5 & 531,8    & \$0,3280 & 0,1360 \\
Gemini 3.8 Flash    & C++        & 53,9 & 1.172,7 & 486,1    & \$0,3065 & 0,1811 \\
Gemini 3.8 Flash    & Java       & 77,7 & 2.209,4 & 668,0    & \$0,3918 & 0,1213 \\
Gemini 3.8 Flash    & JavaScript & 67,9 & 1.658,8 & 565,2    & \$0,3436 & 0,0742 \\
\midrule
GPT-6 Luna          & Python     & 34,2 & 778,5   & 1.361,0  & \$0,0948 & 0,3595 \\
GPT-6 Luna          & C          & 56,4 & 1.316,7 & 1.527,8  & \$0,1053 & 0,1734 \\
GPT-6 Luna          & C++        & 47,3 & 1.071,9 & 1.242,5  & \$0,0874 & 0,2048 \\
GPT-6 Luna          & Java       & 72,3 & 2.066,4 & 1.597,5  & \$0,1096 & 0,1517 \\
GPT-6 Luna          & JavaScript & 46,4 & 1.040,8 & 1.414,2  & \$0,0982 & 0,1326 \\
\bottomrule
\end{tabular}%
}
\end{table}

Os números da Tabela~\ref{tab:res-pp1-eficiencia} mostram como o esforço de geração difere entre os modelos. No \textbf{DeepSeek v4.1 Flash}, a necessidade de construir estruturas de dados manualmente em \linguagem{C} fez a média de saída subir para $11.387{,}3$ \textit{tokens} por questão (contra $7.541{,}5$ em \linguagem{C++}, onde a STL já oferece essas estruturas prontas), resultando em um gasto de US\$\,0,8544 apenas nessa linguagem e US\$\,3,4333 no somatório das cinco baterias. Já o \textbf{Gemini 3.8 Flash} emitiu respostas diretas com média global de $537{,}4$ \textit{tokens} de saída (custo total de US\$\,1,6527), enquanto o \textbf{GPT-6 Luna} apresentou a operação mais econômica da PP1, somando US\$\,0,4953 para avaliar todas as 625 submissões.
\FloatBarrier

\section{Detalhamento Complementar da PP2: Impacto Visual e Modalidades}\label{ap:pp2-detalhes}

\subsection{Contraste entre Questões com e sem Imagens no \textit{Baseline}}\label{sec:res-rq2-baseline}

Para verificar em que medida a presença de figuras no enunciado original afeta os modelos quando submetidos apenas ao texto puro (\modalidade{text}, $T = 0$), a Tabela~\ref{tab:res-pp2-baseline} e a Figura~\ref{fig:res-pp2-img-noimg} comparam o desempenho entre as 46 questões que possuem imagens e as 79 questões puramente textuais. O comparativo é apresentado sob dois recortes: na melhor linguagem individual de cada LLM na PP1 (\linguagem{C++} para o DeepSeek, \linguagem{Python} para o Gemini e \linguagem{Java} para o GPT-6 Luna) e no somatório das cinco linguagens.

\begin{table}[!htbp]
\centering
\caption{Desempenho no \textit{baseline} (\modalidade{text}, $T=0$) entre questões com e sem imagens.}
\label{tab:res-pp2-baseline}
\resizebox{\textwidth}{!}{%
\begin{tabular}{@{}llccccl@{}}
\toprule
\textbf{Recorte} & \textbf{LLM} & \textbf{Ling.} & \textbf{Com Imagens} & \textbf{Sem Imagens} & \textbf{Dif. (p.p.)} & \textbf{Erros (Com Img.)} \\
\midrule
Melhor linguagem  & DeepSeek v4.1 Flash & C++    & $44/46$ ($95{,}65\%$)  & $77/79$ ($97{,}47\%$)  & $-1{,}82$ & $1$ \veredito{WA}, $1$ \veredito{CE} \\
Melhor linguagem  & Gemini 3.8 Flash    & Python & $44/46$ ($95{,}65\%$)  & $77/79$ ($97{,}47\%$)  & $-1{,}82$ & $2$ \veredito{WA} \\
Melhor linguagem  & GPT-6 Luna          & Java   & $42/46$ ($91{,}30\%$)  & $77/79$ ($97{,}47\%$)  & $-6{,}16$ & $4$ \veredito{WA} \\
\midrule
Agregado (5 ling.) & DeepSeek v4.1 Flash & Todas  & $212/230$ ($92{,}17\%$) & $377/395$ ($95{,}44\%$) & $-3{,}27$ & --- \\
Agregado (5 ling.) & Gemini 3.8 Flash    & Todas  & $213/230$ ($92{,}61\%$) & $377/395$ ($95{,}44\%$) & $-2{,}83$ & --- \\
Agregado (5 ling.) & GPT-6 Luna          & Todas  & $201/230$ ($87{,}39\%$) & $365/395$ ($92{,}41\%$) & $-5{,}01$ & --- \\
\bottomrule
\end{tabular}%
}
\end{table}

\begin{figure}[!htbp]
\centering
% Gráfico comparativo: Com Imagens (46) vs Sem Imagens (79) no baseline textual (Cap. 5, Seção 5.3.1)
\begin{tikzpicture}
\begin{groupplot}[
  group style={
    group size=2 by 1,
    horizontal sep=1.3cm
  },
  width=0.47\textwidth,
  height=6.1cm,
  symbolic x coords={DeepSeek,Gemini,GPT6},
  xtick={DeepSeek,Gemini,GPT6},
  enlarge x limits=0.25,
  ymajorgrids=true,
  grid style={black!12},
  axis lines*=left,
  tick label style={font=\scriptsize},
  x tick label style={align=center},
  label style={font=\footnotesize},
  title style={font=\footnotesize\bfseries, yshift=3pt},
  ymin=82, ymax=104.5,
  ytick={85,90,95,100}
]
% Painel (a): Melhor linguagem individual
\nextgroupplot[
  ybar=3pt,
  bar width=14pt,
  title={(a) Melhor linguagem por modelo},
  ylabel={Taxa \passat{1} (\%)},
  xticklabels={DeepSeek\\(C++),Gemini\\(Python),GPT-6\\(Java)},
  legend to name=legpp2imgnoimg,
  legend style={
    legend columns=2,
    draw=black!30,
    rounded corners=2pt,
    font=\footnotesize,
    /tikz/every even column/.append style={column sep=14pt}
  },
  nodes near coords,
  nodes near coords style={
    font=\tiny\bfseries,
    black,
    rotate=90,
    anchor=west,
    /pgf/number format/fixed,
    /pgf/number format/fixed zerofill,
    /pgf/number format/precision=2,
    /pgf/number format/use comma
  }
]
\addplot[fill=vVerde, draw=vVerde!70!black, line width=0.5pt] coordinates {
  (DeepSeek,97.47) (Gemini,97.47) (GPT6,97.47)
};
\addplot[fill=vVermelho, draw=vVermelho!70!black, line width=0.5pt] coordinates {
  (DeepSeek,95.65) (Gemini,95.65) (GPT6,91.30)
};
\legend{Sem Imagens ($n=79$), Com Imagens ($n=46$)}

% Painel (b): Agregado nas 5 linguagens
\nextgroupplot[
  ybar=3pt,
  bar width=14pt,
  title={(b) Agregado nas 5 linguagens},
  ylabel={Taxa média de \veredito{AC} (\%)},
  xticklabels={DeepSeek\\(5 ling.),Gemini\\(5 ling.),GPT-6\\(5 ling.)},
  nodes near coords,
  nodes near coords style={
    font=\tiny\bfseries,
    black,
    rotate=90,
    anchor=west,
    /pgf/number format/fixed,
    /pgf/number format/fixed zerofill,
    /pgf/number format/precision=2,
    /pgf/number format/use comma
  }
]
\addplot[fill=vVerde, draw=vVerde!70!black, line width=0.5pt] coordinates {
  (DeepSeek,95.44) (Gemini,95.44) (GPT6,92.41)
};
\addplot[fill=vVermelho, draw=vVermelho!70!black, line width=0.5pt] coordinates {
  (DeepSeek,92.17) (Gemini,92.61) (GPT6,87.39)
};
\end{groupplot}

% Legenda compartilhada acima dos dois gráficos
\node[anchor=south] at ($(group c1r1.north)!0.5!(group c2r1.north) + (0,0.75cm)$)
  {\pgfplotslegendfromname{legpp2imgnoimg}};
\end{tikzpicture}

\caption[Contraste entre questões com e sem imagens na PP2]{Taxa de aprovação (\passat{1}) no \textit{baseline} textual (\modalidade{text}) para questões com e sem imagens (a) na melhor linguagem de cada LLM e (b) no agregado das cinco linguagens.}
\label{fig:res-pp2-img-noimg}
\end{figure}
\FloatBarrier

Nas 79 questões puramente textuais, os três modelos convergiram para exatamente $77/79$ acertos ($97{,}47\%$) em suas melhores linguagens. Quando avaliados nos 46 problemas cujos cadernos originais trazem figuras --- frequentemente usadas para ilustrar malhas, tabuleiros, cortes geométricos ou ligações em grafos ---, a taxa de aprovação recua em todos os cenários ($-1{,}82\text{ p.p.}$ a $-6{,}16\text{ p.p.}$ na melhor linguagem, e $-2{,}83\text{ p.p.}$ a $-5{,}01\text{ p.p.}$ no agregado das cinco linguagens), confirmando que a ausência da referência visual dificulta a interpretação espacial do enunciado.

\subsection{Transições de Veredito e Telemetria de Custos nas Modalidades Visuais}\label{sec:res-rq2-custo}

No experimento controlado em \linguagem{C++} sobre as 46 questões com imagens (Seção~\ref{sec:res-rq2}), o tempo médio de execução das soluções no \textit{JudgeService} manteve-se praticamente inalterado entre as modalidades \modalidade{text}, \modalidade{img} e \modalidade{description} (variando apenas entre $0{,}164\,\mathrm{s}$ e $0{,}187\,\mathrm{s}$). Das 46 questões avaliadas, 31 ($67{,}39\%$) receberam veredito \veredito{AC} unânime nos três modelos e nas três modalidades. As 15 questões que apresentaram variação de veredito ou falha em ao menos uma configuração estão listadas na Tabela~\ref{tab:res-pp2-divergencias}.

\begin{table}[!htbp]
\centering
\caption{Mapa de transições de vereditos nas $15$ questões com divergência ou falha em \linguagem{C++} (\modalidade{text} / \modalidade{img} / \modalidade{description}).}
\label{tab:res-pp2-divergencias}
\resizebox{\textwidth}{!}{%
\begin{tabular}{@{}lccc@{}}
\toprule
\textbf{Problema} & \textbf{DeepSeek v4.1 Flash} & \textbf{Gemini 3.8 Flash} & \textbf{GPT-6 Luna} \\
\midrule
Arco e flecha          & \veredito{AC} / \veredito{CE} / \veredito{AC} & \veredito{AC} / \veredito{AC} / \veredito{AC} & \veredito{AC} / \veredito{AC} / \veredito{AC} \\
Brigadeiros            & \veredito{AC} / \veredito{AC} / \veredito{WA} & \veredito{AC} / \veredito{AC} / \veredito{WA} & \veredito{WA} / \veredito{WA} / \veredito{WA} \\
Burocracia             & \veredito{AC} / \veredito{AC} / \veredito{AC} & \veredito{AC} / \veredito{AC} / \veredito{AC} & \veredito{AC} / \veredito{WA} / \veredito{WA} \\
Chuva\_2019            & \veredito{CE} / \veredito{AC} / \veredito{AC} & \veredito{AC} / \veredito{AC} / \veredito{AC} & \veredito{WA} / \veredito{AC} / \veredito{AC} \\
Construção de Rodovia  & \veredito{WA} / \veredito{WA} / \veredito{WA} & \veredito{WA} / \veredito{WA} / \veredito{WA} & \veredito{WA} / \veredito{WA} / \veredito{WA} \\
Cortando o Papel       & \veredito{AC} / \veredito{AC} / \veredito{AC} & \veredito{AC} / \veredito{AC} / \veredito{AC} & \veredito{WA} / \veredito{AC} / \veredito{AC} \\
Dario e Xerxes         & \veredito{AC} / \veredito{AC} / \veredito{AC} & \veredito{AC} / \veredito{AC} / \veredito{AC} & \veredito{WA} / \veredito{AC} / \veredito{AC} \\
Fast-Food              & \veredito{AC} / \veredito{AC} / \veredito{AC} & \veredito{WA} / \veredito{AC} / \veredito{AC} & \veredito{AC} / \veredito{AC} / \veredito{AC} \\
Frete                  & \veredito{AC} / \veredito{CE} / \veredito{AC} & \veredito{AC} / \veredito{AC} / \veredito{AC} & \veredito{AC} / \veredito{AC} / \veredito{AC} \\
Metrô da Nlogônia      & \veredito{AC} / \veredito{WA} / \veredito{AC} & \veredito{AC} / \veredito{AC} / \veredito{AC} & \veredito{AC} / \veredito{WA} / \veredito{WA} \\
Nova avenida           & \veredito{AC} / \veredito{AC} / \veredito{WA} & \veredito{AC} / \veredito{AC} / \veredito{AC} & \veredito{AC} / \veredito{AC} / \veredito{AC} \\
Passatempo             & \veredito{AC} / \veredito{AC} / \veredito{AC} & \veredito{AC} / \veredito{AC} / \veredito{AC} & \veredito{CE} / \veredito{WA} / \veredito{AC} \\
Soma                   & \veredito{AC} / \veredito{AC} / \veredito{AC} & \veredito{WA} / \veredito{WA} / \veredito{AC} & \veredito{AC} / \veredito{AC} / \veredito{AC} \\
Transportes            & \veredito{AC} / \veredito{CE} / \veredito{AC} & \veredito{AC} / \veredito{AC} / \veredito{AC} & \veredito{AC} / \veredito{AC} / \veredito{AC} \\
Trio de Bonecas        & \veredito{AC} / \veredito{AC} / \veredito{WA} & \veredito{AC} / \veredito{AC} / \veredito{AC} & \veredito{WA} / \veredito{WA} / \veredito{WA} \\
\bottomrule
\end{tabular}%
}
\end{table}

O mapa individualizado permite enxergar onde o recurso visual realmente ajudou. Questões com forte apelo espacial, como \textit{Chuva\_2019}, \textit{Cortando o Papel}, \textit{Dario e Xerxes} e \textit{Fast-Food}, passaram de erro para aprovação (\veredito{AC}) assim que os modelos receberam a imagem ou sua descrição. Em \textit{Passatempo} (no GPT-6 Luna) e \textit{Soma} (no Gemini 3.8 Flash), apenas a descrição estruturada (\modalidade{description}) foi capaz de conduzir ao acerto. Em contrapartida, no \textbf{DeepSeek v4.1 Flash}, questões que já eram resolvidas apenas com o texto (\textit{Arco e flecha}, \textit{Frete} e \textit{Transportes}) sofreram erro de compilação (\veredito{CE}) quando a imagem foi anexada em Base64, voltando a passar normalmente na modalidade descritiva.

A Tabela~\ref{tab:res-pp2-custos} complementa a Figura~\ref{fig:res-pp2-tokens-custo} do Capítulo~\ref{cap:resultados} com os valores numéricos exatos de \textit{tokens} médios de entrada e saída e o custo total das 46 questões em \linguagem{C++}. Nota-se que substituir a imagem em Base64 pela descrição textual evita o salto de \textit{tokens} de entrada no Gemini ($2.173{,}1 \rightarrow 1.046{,}6$) e o aumento da cadeia de saída no DeepSeek ($12.993{,}4 \rightarrow 5.639{,}5$), reduzindo o custo financeiro sem perda de acurácia.

\begin{table}[!htbp]
\centering
\caption{Consumo médio de \textit{tokens} por questão e custo total em dólares nas $46$ questões em \linguagem{C++} (PP2).}
\label{tab:res-pp2-custos}
\small
\resizebox{\textwidth}{!}{%
\begin{tabular}{@{}llrrr@{}}
\toprule
\textbf{LLM} & \textbf{Modalidade} & \textbf{Tokens Ent. (média)} & \textbf{Tokens Saída (média)} & \textbf{Custo Total (46 q.)} \\
\midrule
DeepSeek v4.1 Flash & \modalidade{text}        & $950{,}8$   & $8.598{,}9$  & \$0,2375 \\
DeepSeek v4.1 Flash & \modalidade{img}         & $1.181{,}3$ & $12.993{,}4$ & \$0,3588 \\
DeepSeek v4.1 Flash & \modalidade{description} & $1.110{,}0$ & $5.639{,}5$  & \$0,1558 \\
\midrule
Gemini 3.8 Flash    & \modalidade{text}        & $901{,}5$   & $454{,}5$    & \$0,1095 \\
Gemini 3.8 Flash    & \modalidade{img}         & $2.173{,}1$ & $482{,}4$    & \$0,1582 \\
Gemini 3.8 Flash    & \modalidade{description} & $1.046{,}6$ & $466{,}5$    & \$0,1166 \\
\midrule
GPT-6 Luna          & \modalidade{text}        & $838{,}2$   & $1.257{,}4$  & \$0,0328 \\
GPT-6 Luna          & \modalidade{img}         & $983{,}0$   & $1.370{,}8$  & \$0,0360 \\
GPT-6 Luna          & \modalidade{description} & $982{,}6$   & $1.125{,}9$  & \$0,0304 \\
\bottomrule
\end{tabular}%
}
\end{table}
\FloatBarrier

\section{Detalhamento Complementar da PP3: Limites de Tempo e Memória em C}\label{ap:pp3-detalhes}

\subsection{Transições de Veredito Questão a Questão}\label{sec:res-rq3-overeng}

A Figura~\ref{fig:res-pp3-erros} ilustra visualmente a migração de falhas entre o \textit{baseline} em \linguagem{C} e a bateria com limites explícitos (PP3). Como programas em \linguagem{C} compilados com \texttt{-O2} já rodam muito abaixo dos tetos da OBI (sem nenhum \veredito{TLE} ou \veredito{MLE}), informar \texttt{time\_limit} e \texttt{memory\_limit} no \textit{prompt} levou os modelos a tentarem otimizações desnecessárias de alocação e leitura, aumentando o total de respostas incorretas (\veredito{WA}).

\begin{figure}[!htbp]
\centering
% Gráfico em 3 painéis empilhados: distribuição de erros (Não-AC) Baseline vs PP3 em C (Cap. 5, Seção 5.4.1)
\begin{tikzpicture}
\begin{groupplot}[
  group style={
    group size=3 by 1,
    horizontal sep=1.1cm,
    yticklabels at=edge left
  },
  width=0.295\textwidth,
  height=5.8cm,
  symbolic x coords={Baseline,PP3},
  xtick={Baseline,PP3},
  xticklabels={Base\\(Sem Lim.),PP3\\(Limites)},
  x tick label style={align=center},
  enlarge x limits=0.45,
  ymajorgrids=true,
  grid style={black!12},
  axis lines*=left,
  tick label style={font=\scriptsize},
  label style={font=\footnotesize},
  title style={font=\footnotesize\bfseries, yshift=2pt},
  ymin=0, ymax=24
]
% Painel (a): DeepSeek v4.1 Flash
\nextgroupplot[
  ybar stacked,
  bar width=22pt,
  title={(a) DeepSeek v4.1 Flash},
  ylabel={Erros (de 125 questões)},
  legend to name=legpp3erros,
  legend style={
    legend columns=3,
    draw=black!30,
    rounded corners=2pt,
    font=\footnotesize,
    /tikz/every even column/.append style={column sep=12pt}
  }
]
\addplot[fill=corWA, draw=corWA!75!black, line width=0.4pt] coordinates {(Baseline,3) (PP3,3)};
\addplot[fill=corCE, draw=corCE!75!black, line width=0.4pt] coordinates {(Baseline,4) (PP3,3)};
\addplot[fill=corRE, draw=corRE!75!black, line width=0.4pt] coordinates {(Baseline,0) (PP3,1)};
\legend{\veredito{WA}, \veredito{CE}, \veredito{RE}}
\node[font=\scriptsize\bfseries, anchor=south] at (axis cs:Baseline,7) {7};
\node[font=\scriptsize\bfseries, anchor=south] at (axis cs:PP3,7) {7};

% Painel (b): Gemini 3.8 Flash
\nextgroupplot[
  ybar stacked,
  bar width=22pt,
  title={(b) Gemini 3.8 Flash}
]
\addplot[fill=corWA, draw=corWA!75!black, line width=0.4pt] coordinates {(Baseline,3) (PP3,6)};
\addplot[fill=corCE, draw=corCE!75!black, line width=0.4pt] coordinates {(Baseline,5) (PP3,3)};
\addplot[fill=corRE, draw=corRE!75!black, line width=0.4pt] coordinates {(Baseline,0) (PP3,0)};
\node[font=\scriptsize\bfseries, anchor=south] at (axis cs:Baseline,8) {8};
\node[font=\scriptsize\bfseries, anchor=south] at (axis cs:PP3,9) {9};

% Painel (c): GPT-6 Luna
\nextgroupplot[
  ybar stacked,
  bar width=22pt,
  title={(c) GPT-6 Luna}
]
\addplot[fill=corWA, draw=corWA!75!black, line width=0.4pt] coordinates {(Baseline,11) (PP3,15)};
\addplot[fill=corCE, draw=corCE!75!black, line width=0.4pt] coordinates {(Baseline,4) (PP3,5)};
\addplot[fill=corRE, draw=corRE!75!black, line width=0.4pt] coordinates {(Baseline,0) (PP3,0)};
\node[font=\scriptsize\bfseries, anchor=south] at (axis cs:Baseline,15) {15};
\node[font=\scriptsize\bfseries, anchor=south] at (axis cs:PP3,20) {20};
\end{groupplot}

% Legenda compartilhada acima dos três painéis
\node[anchor=south] at ($(group c1r1.north)!0.5!(group c3r1.north) + (0,0.65cm)$)
  {\pgfplotslegendfromname{legpp3erros}};
\end{tikzpicture}

\caption[Distribuição de erros em C no Baseline versus PP3]{Decomposição das submissões não aceitas em \linguagem{C} por tipo de veredito (\veredito{WA}, \veredito{CE} e \veredito{RE}) entre o \textit{baseline} e a intervenção PP3 nos três modelos.}
\label{fig:res-pp3-erros}
\end{figure}
\FloatBarrier

Embora o saldo agregado apresentado no Capítulo~\ref{cap:resultados} (Tabela~\ref{tab:res-pp3-comparativo}) indique uma variação líquida de $-6$ aprovações ($345 \rightarrow 339$), a análise individual das submissões mostra uma instabilidade bem maior: exatamente 31 das 125 questões ($24{,}8\%$) mudaram de veredito em pelo menos um LLM, conforme listado na Tabela~\ref{tab:res-pp3-transicoes}. No \textbf{Gemini 3.8 Flash}, por exemplo, o saldo líquido de $-1$ \veredito{AC} esconde 13 trocas de resultado (6 recuperações anuladas por 7 novas falhas), enquanto no \textbf{GPT-6 Luna} ocorreram 21 alterações (8 recuperações contra 13 regressões).

\begin{table}[!htbp]
\centering
\caption{Mapeamento completo das $31$ questões com transição de veredito em \linguagem{C} (\textit{Baseline} $\rightarrow$ \texttt{PP3}).}
\label{tab:res-pp3-transicoes}
\resizebox{\textwidth}{!}{%
\begin{tabular}{@{}lcccccc@{}}
\toprule
\textbf{Problema} & \textbf{DeepSeek (B $\rightarrow$ P)} & \textbf{Gemini (B $\rightarrow$ P)} & \textbf{GPT-6 Luna (B $\rightarrow$ P)} & \textbf{Tempo Lim.} & \textbf{Mem. Lim.} & \textbf{Ano} \\
\midrule
A idade de Dona Mônica & \veredito{AC} $\rightarrow$ \veredito{AC} & \veredito{AC} $\rightarrow$ \veredito{AC} & \veredito{AC} $\rightarrow$ \veredito{WA} & 5,0\,s  & 512\,MB   & 2019 \\
Brigadeiros            & \veredito{AC} $\rightarrow$ \veredito{AC} & \veredito{AC} $\rightarrow$ \veredito{WA} & \veredito{WA} $\rightarrow$ \veredito{WA} & 5,0\,s  & 1.024\,MB & 2024 \\
Burocracia             & \veredito{AC} $\rightarrow$ \veredito{AC} & \veredito{WA} $\rightarrow$ \veredito{AC} & \veredito{WA} $\rightarrow$ \veredito{CE} & 5,0\,s  & 1.024\,MB & 2024 \\
Candidatas             & \veredito{AC} $\rightarrow$ \veredito{AC} & \veredito{CE} $\rightarrow$ \veredito{AC} & \veredito{AC} $\rightarrow$ \veredito{CE} & 15,0\,s & 512\,MB   & 2020 \\
Carrinho               & \veredito{AC} $\rightarrow$ \veredito{AC} & \veredito{AC} $\rightarrow$ \veredito{AC} & \veredito{AC} $\rightarrow$ \veredito{WA} & 5,0\,s  & 256\,MB   & 2017 \\
Ciclovias              & \veredito{AC} $\rightarrow$ \veredito{AC} & \veredito{AC} $\rightarrow$ \veredito{CE} & \veredito{AC} $\rightarrow$ \veredito{CE} & 5,0\,s  & 512\,MB   & 2016 \\
Cortando o Papel       & \veredito{AC} $\rightarrow$ \veredito{AC} & \veredito{AC} $\rightarrow$ \veredito{AC} & \veredito{AC} $\rightarrow$ \veredito{WA} & 3,0\,s  & 256\,MB   & 2017 \\
Dança de Formatura     & \veredito{AC} $\rightarrow$ \veredito{AC} & \veredito{CE} $\rightarrow$ \veredito{AC} & \veredito{AC} $\rightarrow$ \veredito{AC} & 3,0\,s  & 512\,MB   & 2024 \\
Dario e Xerxes         & \veredito{AC} $\rightarrow$ \veredito{WA} & \veredito{AC} $\rightarrow$ \veredito{AC} & \veredito{AC} $\rightarrow$ \veredito{AC} & 2,0\,s  & 256\,MB   & 2017 \\
Dona Formiga           & \veredito{AC} $\rightarrow$ \veredito{AC} & \veredito{AC} $\rightarrow$ \veredito{AC} & \veredito{AC} $\rightarrow$ \veredito{WA} & 1,5\,s  & 256\,MB   & 2020 \\
Energia                & \veredito{AC} $\rightarrow$ \veredito{AC} & \veredito{AC} $\rightarrow$ \veredito{CE} & \veredito{WA} $\rightarrow$ \veredito{AC} & 5,0\,s  & 1.024\,MB & 2025 \\
Escadaria              & \veredito{AC} $\rightarrow$ \veredito{AC} & \veredito{CE} $\rightarrow$ \veredito{AC} & \veredito{AC} $\rightarrow$ \veredito{AC} & 5,0\,s  & 1.024\,MB & 2025 \\
Falha de segurança     & \veredito{AC} $\rightarrow$ \veredito{AC} & \veredito{AC} $\rightarrow$ \veredito{AC} & \veredito{AC} $\rightarrow$ \veredito{WA} & 5,0\,s  & 512\,MB   & 2021 \\
Feira de Artesanato    & \veredito{AC} $\rightarrow$ \veredito{AC} & \veredito{AC} $\rightarrow$ \veredito{CE} & \veredito{AC} $\rightarrow$ \veredito{AC} & 2,0\,s  & 512\,MB   & 2025 \\
Fotos de Relíquias     & \veredito{WA} $\rightarrow$ \veredito{AC} & \veredito{AC} $\rightarrow$ \veredito{AC} & \veredito{WA} $\rightarrow$ \veredito{AC} & 5,0\,s  & 1.024\,MB & 2025 \\
Garamana               & \veredito{AC} $\rightarrow$ \veredito{AC} & \veredito{AC} $\rightarrow$ \veredito{AC} & \veredito{AC} $\rightarrow$ \veredito{WA} & 1,0\,s  & 256\,MB   & 2020 \\
Grupos de Trabalho     & \veredito{AC} $\rightarrow$ \veredito{AC} & \veredito{CE} $\rightarrow$ \veredito{AC} & \veredito{CE} $\rightarrow$ \veredito{AC} & 3,0\,s  & 256\,MB   & 2023 \\
Jogo de Pratos         & \veredito{AC} $\rightarrow$ \veredito{AC} & \veredito{AC} $\rightarrow$ \veredito{WA} & \veredito{AC} $\rightarrow$ \veredito{AC} & 5,0\,s  & 1.024\,MB & 2024 \\
Jogo do Poder          & \veredito{AC} $\rightarrow$ \veredito{AC} & \veredito{AC} $\rightarrow$ \veredito{AC} & \veredito{AC} $\rightarrow$ \veredito{WA} & 5,0\,s  & 512\,MB   & 2024 \\
Lista palíndroma       & \veredito{AC} $\rightarrow$ \veredito{AC} & \veredito{AC} $\rightarrow$ \veredito{AC} & \veredito{CE} $\rightarrow$ \veredito{AC} & 3,0\,s  & 512\,MB   & 2021 \\
Mapa                   & \veredito{AC} $\rightarrow$ \veredito{AC} & \veredito{AC} $\rightarrow$ \veredito{AC} & \veredito{AC} $\rightarrow$ \veredito{WA} & 2,0\,s  & 256\,MB   & 2017 \\
Metrô da Nlogônia      & \veredito{AC} $\rightarrow$ \veredito{AC} & \veredito{AC} $\rightarrow$ \veredito{WA} & \veredito{WA} $\rightarrow$ \veredito{WA} & 5,0\,s  & 1.024\,MB & 2019 \\
Pô, que mão            & \veredito{AC} $\rightarrow$ \veredito{AC} & \veredito{AC} $\rightarrow$ \veredito{AC} & \veredito{CE} $\rightarrow$ \veredito{AC} & 5,0\,s  & 512\,MB   & 2016 \\
Quadrado               & \veredito{WA} $\rightarrow$ \veredito{WA} & \veredito{AC} $\rightarrow$ \veredito{WA} & \veredito{WA} $\rightarrow$ \veredito{WA} & 2,0\,s  & 256\,MB   & 2022 \\
Ralouim                & \veredito{AC} $\rightarrow$ \veredito{AC} & \veredito{AC} $\rightarrow$ \veredito{AC} & \veredito{WA} $\rightarrow$ \veredito{AC} & 10,0\,s & 1.024\,MB & 2020 \\
Soma                   & \veredito{CE} $\rightarrow$ \veredito{AC} & \veredito{AC} $\rightarrow$ \veredito{AC} & \veredito{WA} $\rightarrow$ \veredito{AC} & 5,0\,s  & 1.024\,MB & 2019 \\
Times                  & \veredito{AC} $\rightarrow$ \veredito{AC} & \veredito{CE} $\rightarrow$ \veredito{AC} & \veredito{AC} $\rightarrow$ \veredito{CE} & 5,0\,s  & 512\,MB   & 2016 \\
Torre de dados         & \veredito{AC} $\rightarrow$ \veredito{AC} & \veredito{AC} $\rightarrow$ \veredito{AC} & \veredito{AC} $\rightarrow$ \veredito{WA} & 5,0\,s  & 256\,MB   & 2020 \\
Trem da mina           & \veredito{AC} $\rightarrow$ \veredito{AC} & \veredito{AC} $\rightarrow$ \veredito{AC} & \veredito{AC} $\rightarrow$ \veredito{WA} & 15,0\,s & 512\,MB   & 2020 \\
Trio de Bonecas        & \veredito{CE} $\rightarrow$ \veredito{CE} & \veredito{AC} $\rightarrow$ \veredito{AC} & \veredito{WA} $\rightarrow$ \veredito{AC} & 2,0\,s  & 512\,MB   & 2023 \\
Viagem                 & \veredito{AC} $\rightarrow$ \veredito{RE} & \veredito{AC} $\rightarrow$ \veredito{AC} & \veredito{AC} $\rightarrow$ \veredito{AC} & 5,0\,s  & 512\,MB   & 2022 \\
\bottomrule
\end{tabular}%
}
\par\smallskip{\footnotesize Nas demais $94$ questões ($75{,}2\%$), os três modelos mantiveram vereditos idênticos entre o \textit{baseline} e a PP3.}
\end{table}

\subsection{Distribuição de Tempos de Execução e Custos na PP3}\label{sec:res-rq3-custos}

A Figura~\ref{fig:res-pp3-tempo} mostra por que a média de tempo em CPU subiu ligeiramente na PP3 sem que houvesse mudança real na velocidade típica dos programas: enquanto a média aritmética foi puxada por casos isolados de falha lenta (como o \veredito{RE} de $4{,}454\,\mathrm{s}$ em \textit{Viagem} no DeepSeek e soluções \veredito{WA} próximas de $5{,}0\,\mathrm{s}$ no GPT-6 Luna), a mediana de execução manteve-se praticamente constante nos três modelos (em torno de $0{,}16\,\mathrm{s}$).

\begin{figure}[!htbp]
\centering
% Gráficos lado a lado: Tempo médio de execução e Boxplot (sem outliers) em C (Cap. 5, Seção 5.4.2)
\begin{tikzpicture}
\begin{groupplot}[
  group style={
    group size=2 by 1,
    horizontal sep=1.35cm
  },
  width=0.47\textwidth,
  height=6.2cm,
  ymajorgrids=true,
  grid style={black!12},
  axis lines*=left,
  tick label style={font=\scriptsize},
  label style={font=\footnotesize},
  title style={font=\footnotesize\bfseries, yshift=3pt}
]
% Painel (a): Tempo médio no Judge
\nextgroupplot[
  ybar=3pt,
  bar width=14pt,
  symbolic x coords={DeepSeek,Gemini,GPT6},
  xtick={DeepSeek,Gemini,GPT6},
  xticklabels={DeepSeek\\v4.1 Flash,Gemini\\3.8 Flash,GPT-6\\Luna},
  x tick label style={align=center},
  enlarge x limits=0.24,
  ymin=0, ymax=0.33,
  ylabel={Tempo médio no Judge (s)},
  yticklabel style={/pgf/number format/fixed, /pgf/number format/fixed zerofill, /pgf/number format/precision=2, /pgf/number format/use comma},
  title={(a) Tempo médio (inclui \textit{outliers})},
  legend to name=legpp3tempo,
  legend style={
    legend columns=2,
    draw=black!30,
    rounded corners=2pt,
    font=\footnotesize,
    /tikz/every even column/.append style={column sep=14pt}
  },
  nodes near coords,
  nodes near coords style={
    font=\tiny\bfseries,
    black,
    rotate=90,
    anchor=west,
    /pgf/number format/fixed,
    /pgf/number format/fixed zerofill,
    /pgf/number format/precision=3,
    /pgf/number format/use comma
  }
]
\addplot[fill=vAzul, draw=vAzul!70!black, line width=0.5pt] coordinates {
  (DeepSeek,0.1515) (Gemini,0.1307) (GPT6,0.2531)
};
\addplot[fill=vLaranja, draw=vLaranja!70!black, line width=0.5pt] coordinates {
  (DeepSeek,0.1960) (Gemini,0.1489) (GPT6,0.2375)
};
\legend{Baseline (Sem Limites), PP3 (Com Limites)}

% Painel (b): Boxplot sem outliers (IQR + mediana)
\nextgroupplot[
  boxplot/draw direction=y,
  boxplot/box extend=0.48,
  xmin=0.3, xmax=6.7,
  ymin=0.04, ymax=0.27,
  xtick={1,2,3,4,5,6},
  xticklabels={DS\\Base,DS\\PP3,Gem\\Base,Gem\\PP3,GPT\\Base,GPT\\PP3},
  x tick label style={align=center, font=\tiny},
  ylabel={Distribuição sem \textit{outliers} (s)},
  yticklabel style={/pgf/number format/fixed, /pgf/number format/fixed zerofill, /pgf/number format/precision=2, /pgf/number format/use comma},
  title={(b) Dispersão típica (s/ \textit{outliers})}
]
% DeepSeek Base & PP3
\addplot+[boxplot prepared={
  draw position=1,
  lower whisker=0.0740, lower quartile=0.1263, median=0.1620, upper quartile=0.1701, upper whisker=0.2319
}, fill=vAzul, draw=black!80, line width=0.6pt, solid] coordinates {};
\addplot+[boxplot prepared={
  draw position=2,
  lower whisker=0.0767, lower quartile=0.1275, median=0.1659, upper quartile=0.1777, upper whisker=0.2163
}, fill=vLaranja, draw=black!80, line width=0.6pt, solid] coordinates {};
% Gemini Base & PP3
\addplot+[boxplot prepared={
  draw position=3,
  lower whisker=0.0522, lower quartile=0.1100, median=0.1190, upper quartile=0.1627, upper whisker=0.2285
}, fill=vAzul, draw=black!80, line width=0.6pt, solid] coordinates {};
\addplot+[boxplot prepared={
  draw position=4,
  lower whisker=0.0709, lower quartile=0.1197, median=0.1587, upper quartile=0.1740, upper whisker=0.2151
}, fill=vLaranja, draw=black!80, line width=0.6pt, solid] coordinates {};
% GPT-6 Luna Base & PP3
\addplot+[boxplot prepared={
  draw position=5,
  lower whisker=0.1023, lower quartile=0.1358, median=0.1633, upper quartile=0.1740, upper whisker=0.2166
}, fill=vAzul, draw=black!80, line width=0.6pt, solid] coordinates {};
\addplot+[boxplot prepared={
  draw position=6,
  lower whisker=0.0928, lower quartile=0.1267, median=0.1687, upper quartile=0.1789, upper whisker=0.2433
}, fill=vLaranja, draw=black!80, line width=0.6pt, solid] coordinates {};
\end{groupplot}

% Legenda compartilhada acima dos dois painéis
\node[anchor=south] at ($(group c1r1.north)!0.5!(group c2r1.north) + (0,0.75cm)$)
  {\pgfplotslegendfromname{legpp3tempo}};
\end{tikzpicture}

\caption[Tempos médios e distribuição de execução na PP3]{Análise de tempo no \textit{JudgeService} em \linguagem{C}: (a) tempo médio global (sensível a \textit{outliers}) e (b) distribuição por quartis e mediana sem \textit{outliers}.}
\label{fig:res-pp3-tempo}
\end{figure}
\FloatBarrier

Quanto ao consumo de \textit{tokens} e custo financeiro (Tabela~\ref{tab:res-pp3-custos} e Figura~\ref{fig:res-pp3-tokens-custo}), a inclusão das restrições adicionou cerca de 32 \textit{tokens} ao \textit{prompt} de entrada. No \textbf{DeepSeek v4.1 Flash}, a presença dos limites encurtou o raciocínio de saída em $13{,}3\%$ ($11.387{,}3 \rightarrow 9.870{,}5$ \textit{tokens}), reduzindo o custo por solução aceita de US\$\,0,00724 para US\$\,0,00628. Já no \textbf{GPT-6 Luna}, o código extra gerado para tentar otimizar as soluções elevou os \textit{tokens} de saída em $10{,}0\%$ e, somado à queda de 5 acertos, encareceu o custo por \veredito{AC} em $14{,}6\%$ (US\$\,0,00096 $\rightarrow$ US\$\,0,00110).

\begin{table}[!htbp]
\centering
\caption{Consumo médio de \textit{tokens}, custo total e custo por solução aceita em \linguagem{C} (\textit{Baseline} $\rightarrow$ \texttt{PP3}).}
\label{tab:res-pp3-custos}
\resizebox{\textwidth}{!}{%
\begin{tabular}{@{}lrrrr@{}}
\toprule
\textbf{LLM} & \textbf{Tokens Ent. (média)} & \textbf{Tokens Saída (média)} & \textbf{Custo Total (125 q.)} & \textbf{Custo por AC (US\$/AC)} \\
\midrule
DeepSeek v4.1 Flash & $891{,}9 \rightarrow 923{,}9$ ($+32{,}0$)  & $11.387{,}3 \rightarrow 9.870{,}5$ ($-1.516{,}8$) & \$0,8544 $\rightarrow$ \$0,7406 & \$0,00724 $\rightarrow$ \$0,00628 \\
Gemini 3.8 Flash    & $839{,}0 \rightarrow 873{,}2$ ($+34{,}2$)  & $531{,}8 \rightarrow 508{,}1$ ($-23{,}7$)        & \$0,3280 $\rightarrow$ \$0,3200 & \$0,00280 $\rightarrow$ \$0,00276 \\
GPT-6 Luna          & $781{,}2 \rightarrow 811{,}4$ ($+30{,}1$)  & $1.527{,}8 \rightarrow 1.681{,}0$ ($+153{,}1$)    & \$0,1053 $\rightarrow$ \$0,1152 & \$0,00096 $\rightarrow$ \$0,00110 \\
\bottomrule
\end{tabular}%
}
\end{table}

\begin{figure}[!htbp]
\centering
% Gráficos lado a lado: Tokens médios de saída e Custo por AC em C (Cap. 5, Seção 5.4.2)
\begin{tikzpicture}
\begin{groupplot}[
  group style={
    group size=2 by 1,
    horizontal sep=1.45cm
  },
  width=0.46\textwidth,
  height=6.2cm,
  symbolic x coords={DeepSeek,Gemini,GPT6},
  xtick={DeepSeek,Gemini,GPT6},
  xticklabels={DeepSeek\\v4.1 Flash,Gemini\\3.8 Flash,GPT-6\\Luna},
  x tick label style={align=center},
  enlarge x limits=0.25,
  ymajorgrids=true,
  grid style={black!12},
  axis lines*=left,
  tick label style={font=\scriptsize},
  label style={font=\footnotesize},
  title style={font=\footnotesize\bfseries, yshift=3pt}
]
% Painel (a): Tokens médios de saída
\nextgroupplot[
  ybar=3pt,
  bar width=14pt,
  title={(a) Tokens médios de saída em \linguagem{C}},
  ymin=0, ymax=14800,
  scaled y ticks=false,
  ytick={0,3000,6000,9000,12000},
  ylabel={Média de \textit{tokens} de saída},
  yticklabel style={/pgf/number format/fixed, /pgf/number format/1000 sep={.}},
  legend to name=legpp3tokenscusto,
  legend style={
    legend columns=2,
    draw=black!30,
    rounded corners=2pt,
    font=\footnotesize,
    /tikz/every even column/.append style={column sep=14pt}
  },
  nodes near coords,
  nodes near coords style={
    font=\tiny\bfseries,
    black,
    rotate=90,
    anchor=west,
    /pgf/number format/fixed,
    /pgf/number format/precision=0,
    /pgf/number format/1000 sep={.}
  }
]
\addplot[fill=vAzul, draw=vAzul!70!black, line width=0.5pt] coordinates {
  (DeepSeek,11387.3) (Gemini,531.8) (GPT6,1527.8)
};
\addplot[fill=vLaranja, draw=vLaranja!70!black, line width=0.5pt] coordinates {
  (DeepSeek,9870.5) (Gemini,508.1) (GPT6,1681.0)
};
\legend{Baseline (Sem Limites), PP3 (Com Limites)}

% Painel (b): Custo por solução aceita (USD/AC)
\nextgroupplot[
  ybar=3pt,
  bar width=14pt,
  title={(b) Custo médio por solução aceita},
  ymin=0, ymax=0.0098,
  ylabel={Custo por \veredito{AC} (US\$/\veredito{AC})},
  ytick={0,0.002,0.004,0.006,0.008},
  yticklabel style={/pgf/number format/fixed, /pgf/number format/fixed zerofill, /pgf/number format/precision=3, /pgf/number format/use comma},
  scaled y ticks=false,
  nodes near coords,
  nodes near coords style={
    font=\tiny\bfseries,
    black,
    rotate=90,
    anchor=west,
    /pgf/number format/fixed,
    /pgf/number format/fixed zerofill,
    /pgf/number format/precision=5,
    /pgf/number format/use comma
  }
]
\addplot[fill=vAzul, draw=vAzul!70!black, line width=0.5pt] coordinates {
  (DeepSeek,0.00724) (Gemini,0.00280) (GPT6,0.00096)
};
\addplot[fill=vLaranja, draw=vLaranja!70!black, line width=0.5pt] coordinates {
  (DeepSeek,0.00628) (Gemini,0.00276) (GPT6,0.00110)
};
\end{groupplot}

% Legenda compartilhada acima dos dois gráficos
\node[anchor=south] at ($(group c1r1.north)!0.5!(group c2r1.north) + (0,0.75cm)$)
  {\pgfplotslegendfromname{legpp3tokenscusto}};
\end{tikzpicture}

\caption[Tokens de saída e custo por AC em C na PP3]{Telemetria em \linguagem{C} entre \textit{Baseline} e PP3: (a) média de \textit{tokens} de saída gerados por problema e (b) custo financeiro médio por solução aprovada (US\$/AC).}
\label{fig:res-pp3-tokens-custo}
\end{figure}
\FloatBarrier

\section{Detalhamento Complementar da PP4: Estratégias de \texorpdfstring{\textit{Prompting}}{Prompting} em C}\label{ap:pp4-detalhes}

\subsection{Decomposição de Erros e Desempenho por Fase da OBI}\label{sec:res-rq4-comp}

A Figura~\ref{fig:res-pp4-erros} ilustra a distribuição dos vereditos de erro (\veredito{WA}, \veredito{CE}, \veredito{RE} e \veredito{TLE}) correspondente à Tabela~\ref{tab:res-pp4-comparativo} discutida na Seção~\ref{sec:res-rq4}, evidenciando como o \promptname{Few-Shot} eliminou os erros de compilação no Gemini e no GPT-6 Luna, enquanto o \promptname{CoT Pólya} reduziu os erros de lógica (\veredito{WA}) nos três modelos.

\begin{figure}[!htbp]
\centering
% Gráfico em 3 painéis empilhados: taxonomia de erros (Não-AC) por estratégia em C (Cap. 5, Seção 5.5.1)
\begin{tikzpicture}
\begin{groupplot}[
  group style={
    group size=3 by 1,
    horizontal sep=1.15cm,
    yticklabels at=edge left
  },
  width=0.29\textwidth,
  height=5.9cm,
  symbolic x coords={Zero,Few,CoT},
  xtick={Zero,Few,CoT},
  xticklabels={Zero-\\Shot,Few-\\Shot,CoT\\Pólya},
  x tick label style={align=center},
  enlarge x limits=0.28,
  ymajorgrids=true,
  grid style={black!12},
  axis lines*=left,
  tick label style={font=\scriptsize},
  label style={font=\footnotesize},
  title style={font=\scriptsize\bfseries, yshift=2pt},
  ymin=0, ymax=24.5
]
% Painel (a): DeepSeek v4.1 Flash
\nextgroupplot[
  ybar stacked,
  bar width=16pt,
  title={(a) DeepSeek v4.1 Flash},
  ylabel={Erros (de 125 questões)},
  legend to name=legpp4erros,
  legend style={
    legend columns=4,
    draw=black!30,
    rounded corners=2pt,
    font=\footnotesize,
    /tikz/every even column/.append style={column sep=10pt}
  }
]
\addplot[fill=corWA, draw=corWA!75!black, line width=0.4pt] coordinates {(Zero,3) (Few,5) (CoT,3)};
\addplot[fill=corCE, draw=corCE!75!black, line width=0.4pt] coordinates {(Zero,3) (Few,4) (CoT,13)};
\addplot[fill=corRE, draw=corRE!75!black, line width=0.4pt] coordinates {(Zero,1) (Few,0) (CoT,0)};
\addplot[fill=corTLE, draw=corTLE!75!black, line width=0.4pt] coordinates {(Zero,0) (Few,0) (CoT,0)};
\legend{\veredito{WA}, \veredito{CE}, \veredito{RE}, \veredito{TLE}}
\node[font=\scriptsize\bfseries, anchor=south] at (axis cs:Zero,7) {7};
\node[font=\scriptsize\bfseries, anchor=south] at (axis cs:Few,9) {9};
\node[font=\scriptsize\bfseries, anchor=south] at (axis cs:CoT,16) {16};

% Painel (b): Gemini 3.8 Flash
\nextgroupplot[
  ybar stacked,
  bar width=16pt,
  title={(b) Gemini 3.8 Flash}
]
\addplot[fill=corWA, draw=corWA!75!black, line width=0.4pt] coordinates {(Zero,6) (Few,6) (CoT,4)};
\addplot[fill=corCE, draw=corCE!75!black, line width=0.4pt] coordinates {(Zero,3) (Few,0) (CoT,0)};
\addplot[fill=corRE, draw=corRE!75!black, line width=0.4pt] coordinates {(Zero,0) (Few,0) (CoT,0)};
\addplot[fill=corTLE, draw=corTLE!75!black, line width=0.4pt] coordinates {(Zero,0) (Few,0) (CoT,0)};
\node[font=\scriptsize\bfseries, anchor=south] at (axis cs:Zero,9) {9};
\node[font=\scriptsize\bfseries, anchor=south] at (axis cs:Few,6) {6};
\node[font=\scriptsize\bfseries, anchor=south] at (axis cs:CoT,4) {4};

% Painel (c): GPT-6 Luna
\nextgroupplot[
  ybar stacked,
  bar width=16pt,
  title={(c) GPT-6 Luna}
]
\addplot[fill=corWA, draw=corWA!75!black, line width=0.4pt] coordinates {(Zero,15) (Few,10) (CoT,10)};
\addplot[fill=corCE, draw=corCE!75!black, line width=0.4pt] coordinates {(Zero,5) (Few,0) (CoT,9)};
\addplot[fill=corRE, draw=corRE!75!black, line width=0.4pt] coordinates {(Zero,0) (Few,1) (CoT,0)};
\addplot[fill=corTLE, draw=corTLE!75!black, line width=0.4pt] coordinates {(Zero,0) (Few,1) (CoT,0)};
\node[font=\scriptsize\bfseries, anchor=south] at (axis cs:Zero,20) {20};
\node[font=\scriptsize\bfseries, anchor=south] at (axis cs:Few,12) {12};
\node[font=\scriptsize\bfseries, anchor=south] at (axis cs:CoT,19) {19};
\end{groupplot}

% Legenda compartilhada acima dos três painéis
\node[anchor=south] at ($(group c1r1.north)!0.5!(group c3r1.north) + (0,0.65cm)$)
  {\pgfplotslegendfromname{legpp4erros}};
\end{tikzpicture}

\caption[Decomposição de erros em C por estratégia na PP4]{Distribuição das submissões não aceitas em \linguagem{C} por veredito (\veredito{WA}, \veredito{CE}, \veredito{RE} e \veredito{TLE}) nas três estratégias de \textit{prompting}.}
\label{fig:res-pp4-erros}
\end{figure}
\FloatBarrier

Para avaliar em qual nível de dificuldade das provas cada técnica trouxe maior benefício, a Tabela~\ref{tab:res-pp4-fases} estratifica a taxa média de acerto nas edições rotuladas diretamente como \textit{Fase~1} ($n=27$), \textit{Fase~2} ($n=32$) e \textit{Fase~3} ($n=25$), além de contabilizar o balanço de questões recuperadas e degradadas frente ao \textit{Zero-Shot}. Nota-se que o maior salto ocorreu nas questões intermediárias da \textit{Fase~2}, onde tanto o \promptname{Few-Shot} quanto o \promptname{CoT Pólya} elevaram a média de $89{,}58\%$ para $97{,}92\%$ ($+8{,}34\text{ p.p.}$). No balanço individual, o \textbf{Gemini 3.8 Flash} com \promptname{CoT Pólya} recuperou 5 questões sem sofrer nenhuma regressão ($+5/-0$), e o \textbf{GPT-6 Luna} recuperou 12 questões em ambas as técnicas guiadas.

\begin{table}[!htbp]
\centering
\caption{Taxa média de aprovação (\passat{1}) por fase da OBI Sênior em \linguagem{C} e balanço de transições frente ao \textit{Zero-Shot}.}
\label{tab:res-pp4-fases}
\resizebox{\textwidth}{!}{%
\begin{tabular}{@{}lccc|lcc@{}}
\toprule
\textbf{Fase da OBI} & \textbf{\textit{Zero-Shot}} & \textbf{\textit{Few-Shot}} & \textbf{\textit{CoT Pólya}} & \textbf{LLM (Balanço Rec. / Deg.)} & \textbf{\textit{Few-Shot}} & \textbf{\textit{CoT Pólya}} \\
\midrule
Fase 1 ($n=27$) & $100{,}00\%$ & $100{,}00\%$ & $93{,}83\%$ & DeepSeek v4.1 Flash & $+4 / -6$ ($-2$)  & $+4 / -13$ ($-9$) \\
Fase 2 ($n=32$) & $89{,}58\%$  & \tfirst{97,92\%} & \tfirst{97,92\%} & Gemini 3.8 Flash    & $+4 / -1$ ($+3$)  & \tfirst{$+5 / -0$ ($+5$)} \\
Fase 3 ($n=25$) & $81{,}33\%$  & $80{,}00\%$  & $74{,}67\%$ & GPT-6 Luna          & \tfirst{$+12 / -4$ ($+8$)} & $+12 / -11$ ($+1$) \\
\bottomrule
\end{tabular}%
}
\par\smallskip{\footnotesize Colunas à esquerda consideram as edições rotuladas literalmente como \texttt{Fase 1}, \texttt{Fase 2} e \texttt{Fase 3}; à direita, o número de questões recuperadas (Rec.) e degradadas (Deg.) frente ao \textit{Zero-Shot}.}
\end{table}

\subsection{Auditoria dos Erros de Extração no \textit{CoT Pólya}}\label{sec:res-rq4-parser}

Conforme apontado na Seção~\ref{sec:res-rq4}, o aumento de erros de compilação (\veredito{CE}) no \promptname{CoT Pólya} para o \textbf{DeepSeek v4.1 Flash} (13 casos) e para o \textbf{GPT-6 Luna} (9 casos) motivou uma auditoria direta sobre o campo \texttt{raw\_response} dos arquivos JSON de saída. Como o roteiro de Pólya pede que o modelo escreva quatro seções analíticas, essas respostas frequentemente continham blocos Markdown auxiliares (com rascunhos de entrada/saída ou análise na etapa 4) além do bloco \texttt{\textasciigrave\textasciigrave\textasciigrave c ... \textasciigrave\textasciigrave\textasciigrave}, fazendo o extrator baseado no maior bloco genérico capturar texto explicativo em vez do programa em \linguagem{C}.

A Tabela~\ref{tab:res-pp4-parser} resume o resultado dessa verificação sobre os 22 casos de \veredito{CE}. No \textbf{GPT-6 Luna}, 8 das 9 falhas ($88{,}9\%$) continham um programa \linguagem{C} completo com função \texttt{main()} dentro da resposta bruta; no \textbf{DeepSeek v4.1 Flash}, 6 dos 13 casos ($46{,}2\%$) também traziam código \linguagem{C} completo, enquanto outros 5 retornaram objetos JSON estruturados em texto e 2 estouraram o limite de saída. Se o extrator filtrasse estritamente blocos marcados com a linguagem-alvo (\texttt{\textasciigrave\textasciigrave\textasciigrave c}) contendo \texttt{main()}, o GPT-6 Luna poderia alcançar até $91{,}20\%$ ($114/125$) e o DeepSeek até $92{,}00\%$ ($115/125$).

\begin{table}[!htbp]
\centering
\caption{Auditoria dos $22$ casos de \veredito{CE} sob \promptname{CoT Pólya} em \linguagem{C}: códigos válidos ocultos e impacto potencial.}
\label{tab:res-pp4-parser}
\resizebox{\textwidth}{!}{%
\begin{tabular}{@{}lcccp{5.4cm}cc@{}}
\toprule
\textbf{LLM} & \textbf{Total CE} & \textbf{Cód. \texttt{main()}} & \textbf{JSON / Inc.} & \textbf{Questões com Código C Oculto} & \textbf{\passat{1} Bruto} & \textbf{Teto Ajustado*} \\
\midrule
DeepSeek v4.1 Flash & $13$ & $6$ ($46{,}2\%$) & $5$ JSON, $2$ vazios & \textit{Botas Trocadas}, \textit{Camisetas da Olimpíada}, \textit{Chuva\_2022}, \textit{Construção de Rodovia}, \textit{Fotos de Relíquias}, \textit{Remove Dígitos} & $87{,}20\%$ ($109$) & até $92{,}00\%$ ($115$) \\
GPT-6 Luna          & $9$  & $8$ ($88{,}9\%$) & $1$ truncado         & \textit{Fast-Food}, \textit{Festa Junina}, \textit{Festa olímpica}, \textit{Game-10}, \textit{Pirâmide}, \textit{Quadrado}, \textit{Soma}, \textit{Trem da mina} & $84{,}80\%$ ($106$) & até $91{,}20\%$ ($114$) \\
\bottomrule
\end{tabular}%
}
\par\smallskip{\footnotesize *Teto ajustado calculado somando ao \passat{1} bruto os casos que continham bloco \texttt{\textasciigrave\textasciigrave\textasciigrave c} válido com função \texttt{main()}, caso extraídos corretamente e aprovados pelo juiz.}
\end{table}

\subsection{Tempos de Execução e Comparativo Visual de Custos na PP4}\label{sec:res-rq4-economia}

As Figuras~\ref{fig:res-pp4-tempo} e~\ref{fig:res-pp4-tokens-custo} complementam a análise de eficiência da Seção~\ref{sec:res-rq4}. Na Figura~\ref{fig:res-pp4-tempo}, observa-se que as medianas de tempo de CPU das soluções aprovadas permaneceram praticamente idênticas nas três estratégias ($\approx 0{,}16\,\mathrm{s}$). Já a Figura~\ref{fig:res-pp4-tokens-custo} permite visualizar o contraste econômico reportado na Tabela~\ref{tab:res-pp4-custos}: o \promptname{Few-Shot} manteve a quantidade de \textit{tokens} de saída contida (chegando a reduzi-la em $18{,}4\%$ no GPT-6 Luna), enquanto o \promptname{CoT Pólya} elevou a geração textual no Gemini e no GPT-6 Luna devido à escrita das quatro etapas de Pólya.

\begin{figure}[!htbp]
\centering
% Gráficos lado a lado: Tempo médio global e Boxplot (AC sem outliers) por estratégia em C (Cap. 5, Seção 5.5.3)
\begin{tikzpicture}
\begin{groupplot}[
  group style={
    group size=2 by 1,
    horizontal sep=1.45cm
  },
  width=0.46\textwidth,
  height=6.2cm,
  ymajorgrids=true,
  grid style={black!12},
  axis lines*=left,
  tick label style={font=\scriptsize},
  label style={font=\footnotesize},
  title style={font=\footnotesize\bfseries, yshift=3pt}
]
% Painel (a): Tempo médio global por estratégia e LLM
\nextgroupplot[
  ybar=2pt,
  bar width=10pt,
  symbolic x coords={DeepSeek,Gemini,GPT6},
  xtick={DeepSeek,Gemini,GPT6},
  xticklabels={DeepSeek\\v4.1 Flash,Gemini\\3.8 Flash,GPT-6\\Luna},
  x tick label style={align=center},
  enlarge x limits=0.25,
  ymin=0, ymax=0.33,
  ylabel={Tempo médio no Judge (s)},
  yticklabel style={/pgf/number format/fixed, /pgf/number format/fixed zerofill, /pgf/number format/precision=2, /pgf/number format/use comma},
  title={(a) Tempo médio por modelo},
  legend to name=legpp4tempo,
  legend style={
    legend columns=3,
    draw=black!30,
    rounded corners=2pt,
    font=\footnotesize,
    /tikz/every even column/.append style={column sep=12pt}
  },
  nodes near coords,
  nodes near coords style={
    font=\tiny\bfseries,
    black,
    rotate=90,
    anchor=west,
    /pgf/number format/fixed,
    /pgf/number format/fixed zerofill,
    /pgf/number format/precision=3,
    /pgf/number format/use comma
  }
]
\addplot[fill=vAzul, draw=vAzul!70!black, line width=0.4pt] coordinates {
  (DeepSeek,0.196) (Gemini,0.149) (GPT6,0.238)
};
\addplot[fill=vTeal, draw=vTeal!70!black, line width=0.4pt] coordinates {
  (DeepSeek,0.242) (Gemini,0.165) (GPT6,0.247)
};
\addplot[fill=vLaranja, draw=vLaranja!70!black, line width=0.4pt] coordinates {
  (DeepSeek,0.150) (Gemini,0.162) (GPT6,0.164)
};
\legend{Zero-Shot, Few-Shot, CoT Pólya}

% Painel (b): Boxplot sem outliers das soluções AC (agregado das 3 LLMs)
\nextgroupplot[
  boxplot/draw direction=y,
  boxplot/box extend=0.45,
  xmin=0.4, xmax=3.6,
  ymin=0.04, ymax=0.27,
  xtick={1,2,3},
  xticklabels={Zero-Shot,Few-Shot,CoT Pólya},
  x tick label style={align=center, font=\scriptsize},
  ylabel={Tempo de CPU em \veredito{AC} (s)},
  yticklabel style={/pgf/number format/fixed, /pgf/number format/fixed zerofill, /pgf/number format/precision=2, /pgf/number format/use comma},
  title={(b) Distribuição em \veredito{AC} (agregado)}
]
\addplot+[boxplot prepared={
  draw position=1,
  lower whisker=0.0482, lower quartile=0.1243, median=0.1659, upper quartile=0.1766, upper whisker=0.2433
}, fill=vAzul, draw=black!80, line width=0.6pt, solid] coordinates {};
\addplot+[boxplot prepared={
  draw position=2,
  lower whisker=0.0533, lower quartile=0.1219, median=0.1646, upper quartile=0.1757, upper whisker=0.2454
}, fill=vTeal, draw=black!80, line width=0.6pt, solid] coordinates {};
\addplot+[boxplot prepared={
  draw position=3,
  lower whisker=0.0582, lower quartile=0.1214, median=0.1597, upper quartile=0.1737, upper whisker=0.2513
}, fill=vLaranja, draw=black!80, line width=0.6pt, solid] coordinates {};
\end{groupplot}

% Legenda compartilhada acima dos dois painéis
\node[anchor=south] at ($(group c1r1.north)!0.5!(group c2r1.north) + (0,0.75cm)$)
  {\pgfplotslegendfromname{legpp4tempo}};
\end{tikzpicture}

\caption[Tempos médios e distribuição AC por estratégia na PP4]{Desempenho em CPU no \textit{JudgeService} em \linguagem{C}: (a) tempo médio global por modelo e estratégia e (b) boxplot agregado dos tempos de soluções aprovadas (\veredito{AC}, sem \textit{outliers}).}
\label{fig:res-pp4-tempo}
\end{figure}

\begin{figure}[!htbp]
\centering
% Gráficos lado a lado: Tokens médios de saída e Custo por AC por estratégia em C (Cap. 5, Seção 5.5.3)
\begin{tikzpicture}
\begin{groupplot}[
  group style={
    group size=2 by 1,
    horizontal sep=1.45cm
  },
  width=0.46\textwidth,
  height=6.2cm,
  symbolic x coords={DeepSeek,Gemini,GPT6},
  xtick={DeepSeek,Gemini,GPT6},
  xticklabels={DeepSeek\\v4.1 Flash,Gemini\\3.8 Flash,GPT-6\\Luna},
  x tick label style={align=center},
  enlarge x limits=0.25,
  ymajorgrids=true,
  grid style={black!12},
  axis lines*=left,
  tick label style={font=\scriptsize},
  label style={font=\footnotesize},
  title style={font=\footnotesize\bfseries, yshift=3pt}
]
% Painel (a): Tokens médios de saída
\nextgroupplot[
  ybar=2pt,
  bar width=10pt,
  title={(a) Tokens médios de saída},
  ymin=0, ymax=13200,
  scaled y ticks=false,
  ytick={0,3000,6000,9000,12000},
  ylabel={Média de \textit{tokens} de saída},
  yticklabel style={/pgf/number format/fixed, /pgf/number format/1000 sep={.}},
  legend to name=legpp4tokenscusto,
  legend style={
    legend columns=3,
    draw=black!30,
    rounded corners=2pt,
    font=\footnotesize,
    /tikz/every even column/.append style={column sep=12pt}
  },
  nodes near coords,
  nodes near coords style={
    font=\tiny\bfseries,
    black,
    rotate=90,
    anchor=west,
    /pgf/number format/fixed,
    /pgf/number format/precision=0,
    /pgf/number format/1000 sep={.}
  }
]
\addplot[fill=vAzul, draw=vAzul!70!black, line width=0.4pt] coordinates {
  (DeepSeek,9870.5) (Gemini,508.1) (GPT6,1681.0)
};
\addplot[fill=vTeal, draw=vTeal!70!black, line width=0.4pt] coordinates {
  (DeepSeek,9964.2) (Gemini,537.9) (GPT6,1371.6)
};
\addplot[fill=vLaranja, draw=vLaranja!70!black, line width=0.4pt] coordinates {
  (DeepSeek,9872.7) (Gemini,1726.8) (GPT6,2450.7)
};
\legend{Zero-Shot, Few-Shot, CoT Pólya}

% Painel (b): Custo médio por solução aceita (USD/AC)
\nextgroupplot[
  ybar=2pt,
  bar width=10pt,
  title={(b) Custo médio por solução aceita},
  ymin=0, ymax=0.0102,
  ylabel={Custo por \veredito{AC} (US\$/\veredito{AC})},
  ytick={0,0.002,0.004,0.006,0.008},
  yticklabel style={/pgf/number format/fixed, /pgf/number format/fixed zerofill, /pgf/number format/precision=3, /pgf/number format/use comma},
  scaled y ticks=false,
  nodes near coords,
  nodes near coords style={
    font=\tiny\bfseries,
    black,
    rotate=90,
    anchor=west,
    /pgf/number format/fixed,
    /pgf/number format/fixed zerofill,
    /pgf/number format/precision=5,
    /pgf/number format/use comma
  }
]
\addplot[fill=vAzul, draw=vAzul!70!black, line width=0.4pt] coordinates {
  (DeepSeek,0.00628) (Gemini,0.00276) (GPT6,0.00110)
};
\addplot[fill=vTeal, draw=vTeal!70!black, line width=0.4pt] coordinates {
  (DeepSeek,0.00645) (Gemini,0.00374) (GPT6,0.00097)
};
\addplot[fill=vLaranja, draw=vLaranja!70!black, line width=0.4pt] coordinates {
  (DeepSeek,0.00680) (Gemini,0.00776) (GPT6,0.00160)
};
\end{groupplot}

% Legenda compartilhada acima dos dois gráficos
\node[anchor=south] at ($(group c1r1.north)!0.5!(group c2r1.north) + (0,0.75cm)$)
  {\pgfplotslegendfromname{legpp4tokenscusto}};
\end{tikzpicture}

\caption[Tokens de saída e custo por AC por estratégia na PP4]{Análise econômica das estratégias de \textit{prompting} em \linguagem{C}: (a) média de \textit{tokens} de saída por resposta e (b) custo financeiro médio por solução aceita (US\$/AC).}
\label{fig:res-pp4-tokens-custo}
\end{figure}
\FloatBarrier

\section{Detalhamento Complementar da PP5: Evolução Geracional sob Paridade de Juiz}\label{ap:pp5-detalhes}

\subsection{Desempenho Detalhado por Modelo e Efeito da Paridade de Juiz}\label{sec:res-rq5-pass1}

A Tabela~\ref{tab:res-pp5-consolidado} e a Figura~\ref{fig:res-pp5-pass1-erros} detalham os indicadores individuais dos quatro modelos comparados nas 46 questões em \linguagem{Python} \textit{Zero-Shot}, complementando a síntese intrafamília apresentada na Tabela~\ref{tab:res-pp5-intrafamilia} do Capítulo~\ref{cap:resultados}.

\begin{table}[!htbp]
\centering
\caption[Desempenho consolidado na PP5 sob paridade de juiz]{Desempenho consolidado dos quatro modelos nas $46$ questões em \linguagem{Python} \textit{Zero-Shot} ($T = 0{,}0$) sob o mesmo \texttt{JudgeService}.}
\label{tab:res-pp5-consolidado}
\small
\resizebox{\textwidth}{!}{%
\begin{tabular}{@{}lllcccccccrrr@{}}
\toprule
\textbf{Modelo} & \textbf{Família} & \textbf{Geração} & \textbf{AC} & \textbf{\passat{1}} & \textbf{WA} & \textbf{RE} & \textbf{TLE} & \textbf{CE} & \textbf{Tokens Méd.} & \textbf{Custo Total} & \textbf{Custo / AC} \\
\midrule
Gemini 3.1 Pro   & Gemini & Anterior & \tfirst{45}  & \tfirst{97,83\%}  & $1$  & $0$ & $0$ & $0$ & $6.025{,}5$ & \$2,8602 & \$0,06356 \\
Gemini 3.8 Flash & Gemini & Atual    & \tsecond{44} & \tsecond{95,65\%} & $2$  & $0$ & $0$ & $0$ & $2.940{,}4$ & \$0,0994 & \tsecond{\$0,00226} \\
\midrule
GPT-5.4          & GPT    & Anterior & $31$         & $67{,}39\%$       & $12$ & $2$ & $1$ & $0$ & \tfirst{$1.206{,}4$} & \$0,3013 & \$0,00972 \\
GPT-6 Luna       & GPT    & Atual    & $38$         & $82{,}61\%$       & $6$  & $0$ & $2$ & $0$ & \tsecond{$2.044{,}9$} & \tfirst{\$0,0316} & \tfirst{\$0,00083} \\
\bottomrule
\end{tabular}%
}
\par\smallskip{\footnotesize \tfirst{Vermelho} e \tsecond{azul}: melhor e segundo melhor valor nos indicadores-chave.}
\end{table}

\begin{figure}[!htbp]
\centering
% Gráfico em 2 painéis: Pass@1 (%) e Taxonomia de Falhas (Não-AC) na PP5 (Cap. 5, Seção 5.6.1)
\begin{tikzpicture}
\begin{groupplot}[
  group style={
    group size=2 by 1,
    horizontal sep=1.55cm
  },
  width=0.455\textwidth,
  height=6.2cm,
  ymajorgrids=true,
  grid style={black!12},
  axis lines*=left,
  tick label style={font=\scriptsize},
  label style={font=\footnotesize},
  title style={font=\footnotesize\bfseries, yshift=3pt}
]
% Painel (a): Taxa Pass@1 (%) comparando geração anterior vs. geração atual por família
\nextgroupplot[
  ybar=4pt,
  bar width=18pt,
  clip=false,
  title={(a) Taxa \passat{1} em \linguagem{Python} ($n=46$)},
  symbolic x coords={Gemini,GPT},
  xtick={Gemini,GPT},
  xticklabels={Família Gemini\\{\tiny (3.1 Pro $\rightarrow$ 3.8 Flash)},Família GPT\\{\tiny (5.4 $\rightarrow$ 6 Luna)}},
  x tick label style={align=center},
  enlarge x limits=0.44,
  ymin=55, ymax=112,
  ytick={60,70,80,90,100},
  ylabel={Taxa \passat{1} (\%)},
  legend to name=legpp5pass1,
  legend style={
    legend columns=2,
    draw=black!30,
    rounded corners=2pt,
    font=\scriptsize,
    /tikz/every even column/.append style={column sep=10pt}
  },
  nodes near coords,
  nodes near coords style={
    font=\scriptsize\bfseries,
    black,
    yshift=1pt,
    /pgf/number format/fixed,
    /pgf/number format/fixed zerofill,
    /pgf/number format/precision=2,
    /pgf/number format/use comma
  }
]
\addplot[fill=vCinza, draw=vCinza!75!black, line width=0.5pt] coordinates {
  (Gemini,97.83) (GPT,67.39)
};
\addplot[fill=vTeal, draw=vTeal!75!black, line width=0.5pt] coordinates {
  (Gemini,95.65) (GPT,82.61)
};
\legend{Geração anterior, Geração atual}

% Pílulas com a variação (Delta) acima de cada par
\node[font=\tiny\bfseries, fill=vCinza!14, draw=vCinza!75!black, rounded corners=2pt, inner sep=2.2pt]
  at (axis cs:Gemini,107.0) {$\Delta$: $-2{,}18$ p.p.};
\node[font=\tiny\bfseries, fill=vVerde!14, draw=vVerde!75!black, rounded corners=2pt, inner sep=2.2pt]
  at (axis cs:GPT,93.5) {$\Delta$: $+15{,}22$ p.p.};

% Painel (b): Taxonomia de erros (Não-AC) sob o mesmo JudgeService
\nextgroupplot[
  ybar stacked,
  bar width=18pt,
  title={(b) Taxonomia de falhas (mesmo juiz)},
  symbolic x coords={G31,G38,GPT54,GPT6},
  xtick={G31,G38,GPT54,GPT6},
  xticklabels={Gemini\\3.1 Pro,Gemini\\3.8 Flash,GPT-5.4,GPT-6\\Luna},
  x tick label style={align=center},
  enlarge x limits=0.22,
  ymin=0, ymax=19,
  ytick={0,3,6,9,12,15,18},
  ylabel={Erros (de 46 questões)},
  legend to name=legpp5erros,
  legend style={
    legend columns=3,
    draw=black!30,
    rounded corners=2pt,
    font=\scriptsize,
    /tikz/every even column/.append style={column sep=10pt}
  }
]
\addplot[fill=corWA, draw=corWA!75!black, line width=0.4pt] coordinates {
  (G31,1) (G38,2) (GPT54,12) (GPT6,6)
};
\addplot[fill=corRE, draw=corRE!75!black, line width=0.4pt] coordinates {
  (G31,0) (G38,0) (GPT54,2) (GPT6,0)
};
\addplot[fill=corTLE, draw=corTLE!75!black, line width=0.4pt] coordinates {
  (G31,0) (G38,0) (GPT54,1) (GPT6,2)
};
\legend{\veredito{WA}, \veredito{RE}, \veredito{TLE}}

% Totais acima de cada barra empilhada
\node[font=\scriptsize\bfseries, anchor=south] at (axis cs:G31,1) {1};
\node[font=\scriptsize\bfseries, anchor=south] at (axis cs:G38,2) {2};
\node[font=\scriptsize\bfseries, anchor=south] at (axis cs:GPT54,15) {15};
\node[font=\scriptsize\bfseries, anchor=south] at (axis cs:GPT6,8) {8};
\end{groupplot}

% Legendas individuais acima de cada painel
\node[anchor=south] at ($(group c1r1.north) + (0,0.78cm)$)
  {\pgfplotslegendfromname{legpp5pass1}};
\node[anchor=south] at ($(group c2r1.north) + (0,0.78cm)$)
  {\pgfplotslegendfromname{legpp5erros}};
\end{tikzpicture}

\caption[Pass@1 e taxonomia de falhas na PP5]{Comparação entre gerações nas $46$ questões em \linguagem{Python} \textit{Zero-Shot} ($T=0$): (a) taxa \passat{1} intrafamília com variação em pontos percentuais e (b) composição dos erros sob o mesmo \texttt{JudgeService}.}
\label{fig:res-pp5-pass1-erros}
\end{figure}
\FloatBarrier

A inspeção dos vereditos mostra por que o rejulgamento dos códigos da geração anterior no mesmo \texttt{JudgeService} foi indispensável. Na avaliação original do estudo preliminar, o teto fixo de $3{,}0\,\mathrm{s}$ para \linguagem{Python} havia gerado falsos \veredito{TLE} em problemas cujo limite oficial da OBI é de $5{,}0\,\mathrm{s}$, como \textit{Brigadeiros} (no Gemini 3.1 Pro) e \textit{Transportes} (em ambos os modelos anteriores). De forma semelhante, a comparação de saída sem remoção de espaços ao final de cada linha havia penalizado a questão \textit{Pirâmide} (2022) com falso \veredito{WA}. Com a paridade de juiz, o \textbf{Gemini 3.1 Pro} subiu de 42 para 45 acertos e o \textbf{GPT-5.4} passou de 29 para 31 acertos, garantindo que a evolução reportada para o \textbf{GPT-6 Luna} ($38/46$, com redução de \veredito{WA} de 12 para 6 e zeramento dos \veredito{RE}s) represente um ganho real de geração de código.

Quanto à contabilização financeira na família Google, cabe registrar que o custo retornado pela API para o \textbf{Gemini 3.8 Flash} (US\$\,0,0994 no total, ou US\$\,0,00226 por \veredito{AC}) incide sobre os \textit{tokens} visíveis de entrada e saída, enquanto o total de \textit{tokens} reportado na Tabela~\ref{tab:res-pp5-consolidado} inclui $75.571$ \textit{thinking tokens} internos ao longo das 46 questões. Mesmo que todos esses \textit{tokens} internos fossem tarifados pelo valor cheio de saída (US\$\,3,75 por milhão), o custo subiria para US\$\,0,3828 (US\$\,0,00870 por \veredito{AC}), ainda $7{,}3$ vezes menor que o do \textbf{Gemini 3.1 Pro}.

\subsection{Mapeamento das Questões com Divergência entre Gerações}\label{sec:res-rq5-fases}

Das 46 questões avaliadas na PP5, 30 ($65{,}22\%$) obtiveram aprovação unânime nos quatro modelos. As 16 questões em que ocorreu falha ou mudança de veredito entre gerações são apresentadas na Tabela~\ref{tab:res-pp5-divergentes}.

\begin{table}[!htbp]
\centering
\caption[Questões com falha ou divergência entre os quatro modelos na PP5]{Vereditos nas $16$ questões em que ao menos um dos quatro modelos não obteve \veredito{AC} em \linguagem{Python} \textit{Zero-Shot} ($T = 0{,}0$).}
\label{tab:res-pp5-divergentes}
\small
\resizebox{\textwidth}{!}{%
\begin{tabular}{@{}cllcccc@{}}
\toprule
\textbf{Ano} & \textbf{Questão} & \textbf{Etapa} & \textbf{Gemini 3.1 Pro} & \textbf{Gemini 3.8 Flash} & \textbf{GPT-5.4} & \textbf{GPT-6 Luna} \\
\midrule
2016 & Arco e flecha         & Fase 2 / Estadual & \veredito{AC} & \veredito{WA} & \veredito{RE}  & \veredito{AC}  \\
2016 & Nova avenida          & Fase 1 / Local    & \veredito{AC} & \veredito{AC} & \veredito{WA}  & \veredito{AC}  \\
2017 & Cortando o Papel      & Fase 2 / Estadual & \veredito{AC} & \veredito{AC} & \veredito{WA}  & \veredito{AC}  \\
2017 & Dario e Xerxes        & Fase 2 / Estadual & \veredito{AC} & \veredito{AC} & \veredito{WA}  & \veredito{WA}  \\
2019 & Calçada Imperial      & Fase 1 / Local    & \veredito{AC} & \veredito{AC} & \veredito{WA}  & \veredito{AC}  \\
2019 & Chuva\_2019           & Fase 1 / Local    & \veredito{AC} & \veredito{AC} & \veredito{TLE} & \veredito{WA}  \\
2019 & Metrô da Nlogônia     & Fase 3 / Nacional & \veredito{AC} & \veredito{AC} & \veredito{WA}  & \veredito{AC}  \\
2021 & Dona Minhoca          & Fase 3 / Nacional & \veredito{AC} & \veredito{AC} & \veredito{WA}  & \veredito{WA}  \\
2021 & Retângulo             & Fase 2 / Estadual & \veredito{AC} & \veredito{AC} & \veredito{WA}  & \veredito{AC}  \\
2022 & Construção de Rodovia & Fase 3 / Nacional & \veredito{WA} & \veredito{WA} & \veredito{WA}  & \veredito{WA}  \\
2023 & Fast-Food             & Fase 3 / Nacional & \veredito{AC} & \veredito{AC} & \veredito{WA}  & \veredito{TLE} \\
2023 & Trio de Bonecas       & Fase 3 / Nacional & \veredito{AC} & \veredito{AC} & \veredito{WA}  & \veredito{TLE} \\
2024 & Brigadeiros           & Fase 3 / Nacional & \veredito{AC} & \veredito{AC} & \veredito{WA}  & \veredito{AC}  \\
2024 & Burocracia            & Fase 3 / Nacional & \veredito{AC} & \veredito{AC} & \veredito{RE}  & \veredito{WA}  \\
2024 & Jogo da Vida          & Fase 1 / Local    & \veredito{AC} & \veredito{AC} & \veredito{AC}  & \veredito{WA}  \\
2024 & Retas                 & Fase 3 / Nacional & \veredito{AC} & \veredito{AC} & \veredito{WA}  & \veredito{AC}  \\
\bottomrule
\end{tabular}%
}
\par\smallskip{\footnotesize Nas demais $30$ questões ($65{,}22\%$), houve \veredito{AC} unânime.}
\end{table}

A comparação nominal mostra que a família Google manteve estabilidade quase total, registrando apenas uma regressão em \textit{Arco e flecha} (\veredito{AC} $\rightarrow$ \veredito{WA}). Já na família OpenAI, o \textbf{GPT-6 Luna} recuperou 8 problemas que o \textbf{GPT-5.4} havia errado --- com destaque para questões da Fase~2 (\textit{Arco e flecha}, \textit{Cortando o Papel} e \textit{Retângulo}) e da Fase~3 (\textit{Brigadeiros}, \textit{Metrô da Nlogônia} e \textit{Retas}) --- e apresentou apenas uma regressão pontual (\textit{Jogo da Vida}). A única questão sem solução em nenhum dos quatro modelos foi \textit{Construção de Rodovia} (2022), que admite múltiplos pares de saída válidos sem verificador customizado (\textit{special judge}) na base oficial.
\FloatBarrier
